% Introducao.
Os Large Language Models (LLMs) têm se consolidado como uma tecnologia central para obtenção, interpretação e geração de informações em linguagem natural, exercendo impactos significativos em diferentes áreas da sociedade. Aplicações relevantes desses modelos já podem ser observadas na Educação \cite{IAeChatGPT}, com avanços nos métodos de ensino e seu uso como ferramenta de auxílio pedagógico; na Medicina \cite{ChatGPTInMedicice}, oferecendo suporte à análise de dados clínicos e à consulta de literatura biomédica; e na Economia \cite{EconomieAndLLM}, por meio de análises de mercado e avaliações qualitativas de empresas para fins de investimento.

No contexto da pecuária leiteira, torna-se cada vez mais desejável que produtores agropecuários possam interagir com modelos de linguagem por meio de perguntas em linguagem natural, recebendo respostas coerentes, fundamentadas e contextualizadas sobre temas relacionados à atividade. Uma aplicação prática relevante consiste na geração de respostas baseadas em literatura científica para questões relacionadas à alimentação do rebanho, ao manejo animal e ao aumento da produção de leite. Para que isso seja viável, o LLM deve ser enriquecido com conhecimento especializado do domínio, a partir de artigos científicos, documentos técnicos e bases de dados relacionadas à lactação, nutrição e manejo de bovinos leiteiros.

Nesse cenário, técnicas de adaptação e especialização de modelos tornam-se essenciais. Entre elas, destaca-se a Geração Aumentada por Recuperação (Retrieval-Augmented Generation -- RAG) \cite{RAG}, que permite complementar o conhecimento do modelo com informações recuperadas de uma base externa, reduzindo limitações associadas ao conhecimento paramétrico do modelo original. Além do RAG tradicional, abordagens mais recentes, como o LightRAG e o Graph-based RAG (G-RAG), ampliam as possibilidades de recuperação e organização do conhecimento, permitindo explorar diferentes estratégias de indexação, estruturação e recuperação de informações relevantes ao domínio. Paralelamente, técnicas de \textit{prompt engineering} também podem contribuir para melhorar a qualidade das respostas geradas, ao orientar de forma mais precisa o comportamento do modelo diante das consultas realizadas. Adicionalmente, outra estratégia de especialização a ser investigada neste trabalho é o \textit{fine-tuning}, que permite ajustar o modelo diretamente com dados do domínio, buscando maior aderência terminológica e semântica às particularidades da pecuária leiteira.

A avaliação dessas abordagens pode ser realizada por meio de métricas consolidadas de avaliação textual, como BLEU \cite{papineni-etal-2002-bleu}, ROUGE \cite{lin-2004-rouge}, BERTScore \cite{zhang2020bertscoreevaluatingtextgeneration}, BARTScore \cite{bartscore} e GPTScore \cite{fu2023gptscoreevaluatedesire}, as quais possibilitam mensurar a similaridade entre as respostas geradas pelos modelos e os textos de referência. Dessa forma, torna-se possível comparar diferentes estratégias de adaptação e geração de respostas, avaliando seu desempenho em termos de qualidade textual, aderência ao conteúdo técnico e capacidade de resposta no contexto da pecuária leiteira.

Desse modo, este trabalho propõe a análise da aplicação de LLMs associados a uma base de conhecimento especializada em pecuária leiteira, considerando diferentes estratégias de adaptação, incluindo RAG tradicional, LightRAG, G-RAG, técnicas de \textit{prompt engineering} e \textit{fine-tuning}. Busca-se investigar o potencial dessas abordagens como ferramentas de apoio aos pecuaristas, possibilitando a geração de respostas rápidas, detalhadas e tecnicamente fundamentadas, de forma aderente às demandas do contexto produtivo da atividade leiteira.